\section*{ID6311: Fractional Power Law: I(ell\_φ)}
\paragraph{Metadata.}
PiID: 9937162313017 | RTEID: RTE:P24507.0145:T2.6809 | UUID: a70970fb-fb67-56ba-8750-2f8decb5f8da

\paragraph{Canonical Statement.}
\textbackslash{}[ I(\textbackslash{}ell\_\textbackslash{}varphi) = \textbackslash{}int\_0\textasciicircum{}\{\textbackslash{}infty\} \textbackslash{}!r\textbackslash{},dr \textbackslash{}int\_0\textasciicircum{}\{2\textbackslash{}pi\}\textbackslash{}! d\textbackslash{}theta\textbackslash{}, \textbackslash{}psi\_1\textasciicircum{}*(r,\textbackslash{}theta)\textbackslash{}, \textbackslash{}delta\_\textbackslash{}epsilon(r)\textbackslash{}, \textbackslash{}psi\_2(r,\textbackslash{}theta) \textbackslash{}] with \textbackslash{}[ \textbackslash{}psi\_1(r,\textbackslash{}theta)= r\textasciicircum{}\{\textbackslash{}ell\_\textbackslash{}varphi\} e\textasciicircum{}\{i\textbackslash{}ell\_\textbackslash{}varphi\textbackslash{}theta\} e\textasciicircum{}\{-r\textasciicircum{}2/(2\textbackslash{}sigma\textasciicircum{}2)\}, \textbackslash{}quad \textbackslash{}psi\_2(r,\textbackslash{}theta)= e\textasciicircum{}\{-r\textasciicircum{}2/(2\textbackslash{}sigma\textasciicircum{}2)\}. \textbackslash{}]

\paragraph{Canonical Equation.}
\[
I(\ell_\varphi) = \int_0^{\infty} \!r\,dr \int_0^{2\pi}\! d\theta\, \psi_1^*(r,\theta)\, \delta_\epsilon(r)\, \psi_2(r,\theta)
\]

\paragraph{Dictionary.}
Fractional Power Law: I(ell\_φ) — I(ℓ\_φ) = ∫\_0\textasciicircum{}∞ r dr ∫\_0\textasciicircum{}2π dθ ψ\_1\textasciicircum{}*(r,θ) δ\_ε(r) ψ\_2(r,θ)

\paragraph{Encyclopedia.}
Fractional Power Law: I(ell\_φ) is an integral expression, written I(ℓ\_φ) = ∫\_0\textasciicircum{}∞ r dr ∫\_0\textasciicircum{}2π dθ ψ\_1\textasciicircum{}*(r,θ) δ\_ε(r) ψ\_2(r,θ). It is obtained by integrating over a domain, accumulating contributions across the region of interest. It is built from π, θ, ε, r, defining I(ℓ\_φ). Within the Trawin Zero Point registry it serves as one element of the framework's mathematical narrative.
